\subsection{A first-admission floor closes the projected original-score order}
\label{sec:original-score-projected-quartic}

\begin{theorem}[The point-seed original score has a graph-uniform quartic
reserve]
\label{thm:original-score-projected-quartic}
Retain the exact-real point-seed recurrence, interior restricted optima,
frozen zero-padding, transported centers, and named complete gate of
\cref{prop:consumed-energy-structural-solvency}.  Suppose
\begin{equation*}
 0<q<1,
 \qquad \rho=\tau=\frac q5,
 \qquad \Xi_0=R_0=0.
\end{equation*}
In particular, the initial active face is the singleton $\{s\}$ and its
restricted optimum is interior, equivalently
$q^2(1/d_s-q/5)>0$.  Admissions are only the post-step genuinely violating
rows selected by the named gate; no arbitrary closure rows are included.
Then arbitrary nonnegativity projection is permitted and the \emph{original}
all-coordinate score obeys at every reviewed checkpoint
\begin{equation}
 \Xi_k\leq \frac{1-q^2}{q^4}R_k.
 \label{eq:original-score-projected-quartic}
\end{equation}
Consequently the causal ledger is solvent with
\begin{equation*}
 \lambda_{\rm orig}(q):=\frac{1-q^2}{q^4}.
\end{equation*}
Together with \cref{prop:consumed-energy-quartic-lower}, this closes the
graph-uniform coefficient order at $\Theta(q^{-4})$ for the original score,
without an inactivity or normal--anchor hypothesis.  This remains a scalar
solvency theorem: it supplies no convergence, recovery-horizon, semantic
stopping, or work conclusion.
\end{theorem}

\begin{proof}
First observe that the support-positive part of the original score already
obeys
\begin{equation}
 \Xi_{{\rm supp},k}
 \leq \frac{1-q^2}{q^4}R_k.
 \label{eq:original-score-support-part}
\end{equation}
Indeed, the fixed-face decrement in
\cref{eq:support-aware-quartic-reserve} depends only on the projected step,
not on which safe envelope is subsequently used by the gate.  Under any named
frozen admission, every old positive coordinate and its residual row are
unchanged, every new coordinate is zero, and every newly admitted residual
is strictly negative.  Hence $\Xi_{\rm supp}$ is unchanged across the
original gate as well, and the usual induction from $\Xi_0=R_0=0$ proves
\cref{eq:original-score-support-part} along the original chronology.

Before the first admission the face is the singleton $\{s\}$.  If its
projected candidate is positive, the original and support-positive scores
coincide.  If it is zero, then the seed row has negative post-residual by
\cref{eq:clipped-row-post-residual-bound}, so both positive scores are zero.
Thus \cref{eq:original-score-projected-quartic} already holds at every such
checkpoint.  If no admission ever occurs, this finishes the proof.

Suppose instead that the first admission occurs, and choose any admitted
boundary vertex $j$.  Put
\begin{equation*}
 a_q:=\frac{1+q^2}{2},
 \qquad \eta_q:=\frac{1-q^2}{2},
 \qquad x:=\frac1{d_s},
 \qquad K_q:=\frac{2a_q}{5\eta_q}.
\end{equation*}
On the singleton face the load and restricted optimum are
\begin{equation*}
 h_s=q^2\left(x-\frac q5\right)>0,
 \qquad z_s^\star=\frac{h_s}{a_q}.
\end{equation*}
The original and support-aware envelopes coincide on this face, so their
subsolution safety gives $\ell_s\leq z_s^\star$.  The strict admission test is
\begin{equation*}
 \frac{q^3}{5}-\frac{\eta_q}{d_j}\ell_s< -\frac{q^3}{5}.
\end{equation*}
It follows that
\begin{equation}
 \ell_s>\frac{2q^3d_j}{5\eta_q},
 \qquad
 x>q\left(\frac15+K_qd_j\right).
 \label{eq:first-admission-seed-floor}
\end{equation}

Let $E_0$ be the actual transported energy at the zero initialization.  Its
singleton objective-gap component alone gives
\begin{equation*}
 E_0\geq \frac{d_sh_s^2}{2a_q}
 =\frac{q^4}{1+q^2}\frac{(x-q/5)^2}{x}.
\end{equation*}
The scalar function $(x-q/5)^2/x$ is strictly increasing for $x>q/5$.
Using \cref{eq:first-admission-seed-floor}, then minimizing the resulting
bound over the integer $d_j\geq1$, yields
\begin{align}
 E_0
 &>c_0(q)q^5,
 &c_0(q)
 &:=\frac{K_q^2}{(1+q^2)(1/5+K_q)}
 =\frac{4(1+q^2)}{5(1-q^2)(3+q^2)}.
 \label{eq:first-admission-energy-floor}
\end{align}
At the post-step checkpoint that triggers the first admission,
\cref{eq:consumed-reserve-fixed-step} and $R+E=E_0$ on the singleton face
give
\begin{equation}
 R\geq qE_0>c_0(q)q^6.
 \label{eq:first-admission-reserve-floor}
\end{equation}
Exact admission leaves $R$ unchanged, and every later fixed step only
increases it, so this floor persists forever.

It remains to charge coordinates that the support-positive score omits.  If
$p_i=0$ and the post-residual is positive, then $i$ cannot be the seed and
\cref{eq:clipped-row-post-residual-bound} gives $r_i(p)\leq q^3/5$.
Therefore the clipped-zero contribution to the original score is at most
$q^2/25$.  On the other hand,
\begin{align*}
 \frac{1-q^2}{q^4}R
 &>c_0(q)(1-q^2)q^2
 =\frac{4(1+q^2)}{5(3+q^2)}q^2
 >\frac{q^2}{25}.
\end{align*}
The maximum defining the original score is attained either on the positive
support, where \cref{eq:original-score-support-part} applies, or on a clipped
zero row, where the last display applies.  This proves
\cref{eq:original-score-projected-quartic} at all remaining checkpoints.
\end{proof}
